\subsection{\cite{charness2002understanding}'s unilateral dictator games}
\label{sec:cr_app}

CR study a set of simple allocation decisions in which one player (the dictator) chooses between two ways of splitting money with a passive recipient. 
In one version of these games, for example, the dictator (Person B) unilaterally decides between the options ``Left'' and ``Right'':
\begin{quote}
\centering
$\underbrace{(\underbrace{400}_{\mbox{To A}}\:,\: \underbrace{600}_{\mbox{To B}})}_{\mbox{``Left''}}$ $\quad$ vs. $\quad$ $\underbrace{(\underbrace{700}_{\mbox{To A}}\:,\: \underbrace{300}_{\mbox{To B}})}_{\mbox{``Right''}}$.
\end{quote}
CR collected human responses for six variations of this basic dictator setting, each featuring different payoff distributions.
These six settings constitute our training dataset from which we derive the joint empirical distribution $P$ of choosing Left.

Figure~\ref{fig:cr_train} shows the original results from CR.
The columns represent different settings and show the payoffs for each player depending on the dictator's choice of ``Left'' or ``Right''.
The y-axis shows the proportion of the sample that chose ``Left'' for each setting, and the black bars correspond to the distribution of human responses from CR.
Besides the Pareto-dominated Berk23 setting, where everyone chooses ``Right,'' the human data is balanced across the two options.

\begin{figure}[h!]
\begin{center}
\begin{minipage}{\textwidth}
\caption{Distribution of responses for the single-stage training dictator games} 
\label{fig:cr_train}
\vspace*{-0.15in}
\centering
\includegraphics[width=.9\textwidth]{plots/cr_unilater_dict.pdf} 
\end{minipage}
\end{center}
\begin{footnotesize}
\begin{singlespace}
\vspace*{-0.05in}
\emph{Notes:}
This figure reports the results of replications of the unilateral dictator games from \cite{charness2002understanding}. 
Columns each represent a different game, and the x-axis corresponds to different samples of subjects playing each game.
The y-axis shows the proportion of that sample choosing the option ``Left.'' The black bars (and the dashed black lines) are the human responses from the original paper, red is the baseline \aisub{s}, blue are the agents optimized using efficiency, self-interest, inequity aversion as parameters, and the yellow are atheoretical agents with preferences for the TV show new girl, taxidermy, and swimming. 
Error bars report 95\% Wilson confidence intervals.
\end{singlespace}
\end{footnotesize}
\end{figure}

To establish the baseline ($\hat P_0$), we elicited 1,000 responses per setting from \textsc{Gpt-4o}, without any additional instructions.\footnote{\cite{horton2023large} also explore the baseline for the same games. Although their goal is to provide an early demonstration of AI simulation more generally.}
The red bars in Figure~\ref{fig:cr_train} represent these baseline AI responses. Notably, the baseline AI strongly favors choosing ``Right'' in nearly every setting, diverging sharply from the balanced human distributions. 
Quantitatively, this mismatch is substantial: using mean absolute error (MAE) as our distance metric, we find $\frac{1}{6}\sum_{s \in S} d(P_s,\hat P_0)=\CRBaselineTrainMAE$. 
Given that the maximum possible MAE is 1, this is poor baseline predictive accuracy.
